\section{The inverse-permutation condition}
Consider the following condition, as reproduced by Vester and Ham from
Dasgupta and Maskin's definition \citep{vester2012,ham2021,dm1986}:
\begin{equation}\label{eq:wrong}
 u_i(s)=u_{\pi(i)}(P_{\pi^{-1}}s)
       =u_{\pi(i)}(s_{\pi(1)},\ldots,s_{\pi(n)})
 \quad(i\in I,\ s\in A^n,\ \pi\in\Sn).
\end{equation}
Here player $\pi(i)$ receives the action $s_{\pi^2(i)}$, rather than
necessarily receiving $s_i$. The issue is a mismatch of the player index and
profile action; choosing a right-action convention consistently would cause
no error. The original printed definition and its theorem dependencies remain
on our primary-source inspection list.

\begin{theorem}[Vester--Ham characterization]\label{thm:dm}
Let $n\ge3$. A game with a common nonempty action set satisfies
\eqref{eq:wrong} if and only if it is totally symmetric and has identical
payoff functions.
\end{theorem}
\begin{proof}
Suppose \eqref{eq:wrong} holds. For transpositions it agrees with
\eqref{eq:correct}. Transpositions generate $\Sn$, so
Lemma~\ref{lem:group} gives total symmetry.
Comparing \eqref{eq:wrong} for $\pi$ with \eqref{eq:correct} for $\pi^{-1}$
yields
\[
 u_{\pi(i)}(P_{\pi^{-1}}s)
 =u_{\pi^{-1}(i)}(P_{\pi^{-1}}s).
\]
Since $P_{\pi^{-1}}$ is onto, $u_{\pi(i)}=u_{\pi^{-1}(i)}$.
For any distinct $a,b$, choose $k\notin\{a,b\}$ and
$\pi=(k\,a\,b)$. Taking $i=k$ gives $u_a=u_b$.
Conversely, total symmetry applied to $\pi^{-1}$ gives
$u_i(s)=u_{\pi^{-1}(i)}(P_{\pi^{-1}}s)$, which is
\eqref{eq:wrong} when all payoff functions coincide.
\end{proof}
The common-payoff consequence is Vester's Theorem~32; the displayed equivalence
is Ham's Theorem~3.21 \citep{vester2012,ham2021}. The proof above replaces
Ham's auxiliary Lemma~3.20. It does not establish a new priority claim.

\begin{example}[Payoffs need not be constant across profiles]
For three players with $A=\{0,1\}$, put
\[
 u_i(s)=5+7\,\mathbf 1\{s_1+s_2+s_3\ge2\}\quad(i=1,2,3).
\]
The common payoff depends only on the number of ones, so the game is totally
symmetric and satisfies \eqref{eq:wrong}. All payoffs at $(0,0,0)$ are $5$;
all payoffs at $(1,1,0)$ are $12$.
\end{example}

\begin{example}[Sharp distinctions]
For $A=\{0,1\}$, $u_i(s)=s_i$ is totally symmetric. If $n=2$ it also
satisfies \eqref{eq:wrong}, although payoffs differ at $(0,1)$.
For $n\ge3$, choose $\pi=(123)$, $i=1$ and $s=(1,0,0,\ldots,0)$;
then $u_1(s)=1$ while $u_2(P_{\pi^{-1}}s)=0$.
Conversely, $u_i(s)=s_1$ for every $i$ has identical payoff functions but
fails total symmetry. Thus neither part of the characterization may be omitted.
\end{example}

\begin{remark}
For two players every permutation is its inverse, so the two definitions
coincide. For unequal strategy sets raw coordinate permutation may be
undefined. Fixed bijections $b_i:A_i\to A$ permit transport to common
coordinates; arbitrary unrelated pairwise bijections do not specify an action.
The research notes treat this distinction separately.
\end{remark}
